\section{Introducción}
Cuando aumenta la productividad de un sector, suele suponerse que sus empresas están produciendo mejor. Sin embargo, el mismo resultado puede aparecer si las empresas más productivas concentran una mayor proporción del empleo o si dejan de observarse productores de baja productividad. Distinguir estas fuentes cambia el diagnóstico: mejorar las capacidades de las empresas existentes plantea desafíos diferentes a facilitar el crecimiento de las más productivas.

Este informe examina de dónde provienen los cambios de productividad en la industria, el comercio y los servicios colombianos. Utiliza las encuestas anuales empresariales del DANE y la descomposición de Melitz--Polanec para separar cuatro contribuciones: la mejora media de las empresas que permanecen, la asociación entre productividad y participación en el empleo, la entrada y la salida de empresas de la muestra.

Los resultados muestran que no existe una única historia sectorial. En industria, el aumento del índice convive con una contribución interna negativa y depende principalmente de los cambios de composición. En servicios, el periodo reciente incorpora una mejora interna importante. En comercio, el diagnóstico cambia con la ventana analizada y exige especial atención a la continuidad de los registros.

La literatura permite poner estos hallazgos en perspectiva. Encontramos antecedentes de reasignación favorable en Colombia y un patrón de signos semejante en Chile, pero también estudios donde domina la mejora interna. Estas diferencias ayudan a formular preguntas sobre cobertura, horizonte y medición, en lugar de interpretar cualquier coincidencia como una validación.

El informe presenta primero la evidencia nacional e internacional, explica después los datos y el método y desarrolla los resultados sectoriales. El cierre distingue sus implicaciones de las comprobaciones pendientes. Su contribución es un diagnóstico descriptivo: no identifica el efecto causal de una reforma ni mide los costos sociales de la salida empresarial.
\clearpage
